% !PRINTED 116
\weberpage{116}
%<<CONT
und wenn alle Elemente vom $1^{\mathrm{sten}}$ oder $2^{\mathrm{ten}}$ Grade sind, die Gruppen
\[
1, \alpha, \beta, \alpha\beta
\]
mit der Bedingung $\alpha\beta = \beta\alpha$. Es giebt also nur diese zwei Gruppen vom $4^{\mathrm{ten}}$ Grade. Wenn die Elemente mit $1, \alpha, \beta, \gamma$ bezeichnet werden, so haben wir die beiden Tabellen:
\[
\begin{array}{c|cccc}
 & 1 & \alpha & \beta & \gamma \\
\hline
1 & 1 & \alpha & \beta & \gamma \\
\alpha & \alpha & \beta & \gamma & 1 \\
\beta & \beta & \gamma & 1 & \alpha \\
\gamma & \gamma & 1 & \alpha & \beta
\end{array}
\qquad
\begin{array}{c|cccc}
 & 1 & \alpha & \beta & \gamma \\
\hline
1 & 1 & \alpha & \beta & \gamma \\
\alpha & \alpha & 1 & \gamma & \beta \\
\beta & \beta & \gamma & 1 & \alpha \\
\gamma & \gamma & \beta & \alpha & 1
\end{array}
\]

Der Fall $n = 6$ lässt sich in folgender Weise vollständig erledigen: Wir bezeichnen die sechs Elemente mit $1, \alpha, \beta, \gamma, \delta, \varepsilon$, so dass 1 die Einheit der Gruppe ist. Wenn darunter ein Element vom Grade 6 vorkommt, so ist die Gruppe cyklisch.

Wenn wir also von diesem Falle absehen, so können die Grade der Elemente $\alpha, \beta, \gamma, \delta, \varepsilon$ nur 2 oder 3 sein. Sind $\alpha$ und $\beta$ vom $2^{\mathrm{ten}}$ Grade, so kann $\alpha\beta$ nicht vom $2^{\mathrm{ten}}$ Grade sein; denn sonst wäre
\[
\alpha = \alpha^{-1}, \quad \beta = \beta^{-1}, \quad \alpha\beta = \beta^{-1}\alpha^{-1} = \beta\alpha,
\]
und es wäre also $1, \alpha, \beta, \alpha\beta$ eine Gruppe $4^{\mathrm{ten}}$ Grades; eine solche kann aber nicht Theiler einer Gruppe $6^{\mathrm{ten}}$ Grades sein.

Es muss also mindestens ein Element $3^{\mathrm{ten}}$ Grades vorkommen, und wenn wir ein solches mit $\alpha$ bezeichnen und $\gamma$ nicht in $1, \alpha, \alpha^2$ enthalten ist, so ordnet sich die Gruppe so:
\weq{1}{S = 1, \alpha, \alpha^2, \gamma, \gamma\alpha, \gamma\alpha^2.}
Um die Bedingung zu ermitteln, dass dies eine Gruppe sei, bilden wir
\[
\gamma S = \gamma, \gamma\alpha, \gamma\alpha^2, \gamma^2, \gamma^2\alpha, \gamma^2\alpha^2,
\]
was mit $S$ identisch sein muss; und es muss also
\[
\gamma^2 = 1 \text{ oder } = \alpha \text{ oder } = \alpha^2
\]
sein. Ist aber $\gamma^2 = \alpha$ oder $= \alpha^2$, so ist $\gamma^3 = \alpha\gamma$ oder $= \alpha^2\gamma$, und $\gamma^6$, aber keine niedrigere Potenz von $\gamma = 1$, d.\,h. $S$ ist
%CONT>>
